% easychair.tex,v 3.5 2017/03/15

\documentclass{easychair}
%\documentclass[EPiC]{easychair}
%\documentclass[EPiCempty]{easychair}
%\documentclass[debug]{easychair}
%\documentclass[verbose]{easychair}
%\documentclass[notimes]{easychair}
%\documentclass[withtimes]{easychair}
%\documentclass[a4paper]{easychair}
%\documentclass[letterpaper]{easychair}

\usepackage{doc}
\usepackage{minted}
%% \usepackage{amsfonts}
\usepackage{graphicx}
\usepackage{amsmath}
\usepackage{amssymb}
\usepackage{bussproofs}
\usepackage{cite}

% For ≞ (requires the LuaLaTeX engine)
\usepackage{unicode-math}
\setmathfont{Stix Two Math}

% Use this if you have a long article and want to create an index
% \usepackage{makeidx}

% In order to save space or manage large tables or figures in a
% landcape-like text, you can use the rotating and pdflscape
% packages. Uncomment the desired from the below.
%
% \usepackage{rotating}
% \usepackage{pdflscape}

% Some of our commands for this guide.
%
\newcommand{\easychair}{\textsf{easychair}}
\newcommand{\miktex}{MiK{\TeX}}
\newcommand{\texniccenter}{{\TeX}nicCenter}
\newcommand{\makefile}{\texttt{Makefile}}
\newcommand{\latexeditor}{LEd}
\newcommand{\Theory}{\texttt{Theory}}
\newcommand{\Proof}{\texttt{Proof}}
\newcommand{\Proposition}{\texttt{Proposition}}
\newcommand{\Bool}{\texttt{Bool}}
\newcommand{\arrow}{\to}
\newcommand{\limp}{\Rightarrow}
\newcommand{\True}{\texttt{True}}
\newcommand{\False}{\texttt{False}}
\newcommand{\STV}[2]{<\!#1, #2\!>}

%\makeindex

%% Front Matter
%%
% Regular title as in the article class.
%
%% \title{Estimating the Probability of a Proposition to be Provable
%%   Using Probabilistic Logic Networks}
%% \title{Estimating the Probability of a Conjecture to be a Theorem with
%%   Probabilistic Logic Networks, and How to Take Advantage of it for
%%   Inference Control}
\title{Using Forward Chaining to Go Backward}

% Authors are joined by \and. Their affiliations are given by \inst, which indexes
% into the list defined using \institute
%
\author{Nil Geisweiller}

% Institutes for affiliations are also joined by \and,
\institute{
  SingularityNET Foundation,\\
  Zug, Switzerland\\
  \email{nil@singularitynet.io}
}

%  \authorrunning{} has to be set for the shorter version of the authors' names;
% otherwise a warning will be rendered in the running heads. When processed by
% EasyChair, this command is mandatory: a document without \authorrunning
% will be rejected by EasyChair
\authorrunning{Geisweiller}

% \titlerunning{} has to be set to either the main title or its shorter
% version for the running heads. When processed by
% EasyChair, this command is mandatory: a document without \titlerunning
% will be rejected by EasyChair

\titlerunning{Using Forward Chaining to Go Backward}

\begin{document}

\maketitle

%% \begin{abstract}
%% \end{abstract}

\section{Introduction}

If all you have is a forward chainer, can you use it to derive a
process that is functionally equivalent to backward chaining?  This is
the question we attempt to answer in this paper.  This question stems
from the observation that some rewriting systems, such as
MORK~\cite{NEXT}, perform considerably better on forward chaining than
on backward chaining.  In fact, in~\cite{NEXT} we have found that
forwardly generating in MORK the set of all proofs of all theorems up
to a certain depth, then querying the results given any theorem within
that set, is faster than running a naive backward chainer
implementation in MORK.  It might be because our particular
implementation is too naive not well tailored for MORK.  Or it might
be because MORK is inherently bad at backward chaining.  This is
nonetheless to be contrasted with some other rewriting systems such as
PeTTa~\cite{NEXT}, a Prolog-based rewriting system, which seem to
perform equally well on forward and backward chaining~\cite{NEXT}.

Since forward chaining in MORK outperforms forward chaining in PeTTa,
it gives us an incentive to attempt to emulate backward chaining using
forward chaining.  The work presented here target specifically MORK
for that reason, although the idea likely applies to any rewriting
system.

\section{Methodology}

To achieve this feat we decompose the backward chaining task into
three steps
\begin{enumerate}
\item Use forward chaining to generate a theory of backward expansion
  up to some depth.
\item Given some target theorem, apply forward chaining on the theory
  of backward expansion to generate composite rules, such that, if
  they were to run forward, they would derive the given theorem only.
\item Use forward chaining on the obtained composite rules to generate
  proofs of the target theorem.
\end{enumerate}
You may have noticed that these three steps only make use of forward
chaining.

\newpage

%------------------------------------------------------------------------------

\bibliographystyle{plain}
\label{sect:bib}
%\bibliographystyle{alpha}
%\bibliographystyle{unsrt}
%\bibliographystyle{abbrv}
\bibliography{ForwardBackward}

\newpage

%------------------------------------------------------------------------------

\appendix

%------------------------------------------------------------------------------
\end{document}
